\clearpage\begingroup
\SourceLeaf{200}\HistoricalHead{195}{R}{WECHSELWIRKUNGEN ZWISCHEN GALVANISCHEN STRÖMEN UND MAGNETISMUS.}
\begin{center}\textbf{BEMERKUNG.}\end{center}
\noindent\hspace*{\parindent}Die im Vorstehenden abgedruckte Handschrift bildet in inhaltlicher und methodischer Hinsicht eine unmittelbare Fortsetzung und Verallgemeinerung der Zurückführung der Intensität des Erdmagnetismus auf absolutes Maß*).\begingroup\renewcommand{\thefootnote}{}\footnotetext{*) \textit{Intensitas vis magneticae etc.}, Werke V, S. 79.}\endgroup\ Mit den im V. Bande der Werke abgedruckten Stücken zur Elektrodynamik hängt sie nur lose zusammen, höchstens insofern, als einige der dort angegebenen Rechnungen gewissermaßen als Material für die vorliegende Handschrift betrachtet werden könne.

Die Absicht, eine Abhandlung über »die Grundgesetze der Galvanischen Ströme und der Induktion und deren Zurückführung auf absolute Maße« zu schreiben, hat \textsc{Gauss} in einem Briefe an \textsc{Encke} vom Januar 1836 (S. 103 dieses Bandes) ausgesprochen. Er deutet dort auch an, daß »bereits einiges niedergeschrieben« sei. Als Abfassungszeit der vorliegenden Handschrift ist also etwa das Jahr 1835 zu betrachten. Diese Vermutung läßt sich durch folgende Tatsachen stützen: Im Artikel 8. der Handschrift (S. 186) wird von dem Multiplikator im Magnetischen Observatorium gesagt, daß er $2\mathbin{.}135=270$ Windungen Draht besitze. Diese Angabe stimmt mit der Beschreibung des Multiplikators überein, die \textsc{Gauss} in einem vom 4. Dezember 1834 datierten Briefe an \textsc{Gerling} (S. 96 dieses Bandes) gegeben hat. Andrerseits schreibt \textsc{Gauss} an \textsc{Olbers} am 23. Juli 1836 (S. 110 dieses Bandes), daß »der Multiplikator \textls[120]{jetzt} aus 610 Umwindungen übersponnenen Drahtes besteht«. Innerhalb dieser $1\frac12$ Jahre ist also die Abfassung wohl mit Sicherheit zu setzen. Genauer läßt sich der Abfassungstermin folgendermaßen bestimmen: Im Artikel 10. der Handschrift (S. 193) wird der Versuch beschrieben, die magnetische Wirkung der Reibungselektrizität am Multiplikator zu messen. In dem schon genannten Briefe an \textsc{Encke} vom Januar 1836 gibt nun \textsc{Gauss} an, daß er die magnetische Wirkung der Reibungselektrizität »im Laufe des verflossenen Jahres« gemessen habe. Danach ist als die wahrscheinlichste Abfassungszeit das Ende des Jahres 1835 anzusehen.
\begin{flushright}\textsc{Schaefer.}\end{flushright}
\begin{center}\rule{.24\textwidth}{.35pt}\end{center}
\clearpage\endgroup
